\subsection{Krull 整环上的多项式环}

命题 13. 设 A 是 Krull 整环，\(\mathbf{X}_1, \mathbf{X}_2, \ldots, \mathbf{X}_n\) 是不定元。环 \(\mathrm{A}[\mathbf{X}_1, \ldots, \mathbf{X}_n]\) 是 Krull 整环。

对 n 作归纳，只需证明：若 X 是不定元，则 A{[}X{]} 是 Krull 整环。令 K 为 A 的分式域。A{[}X{]} 的分式域是 K(X)。令 I 为 K{[}X{]} 中在 K 上不可约的首一多项式集合；对于所有 \(f \in I\)，令 \(v_{f}\) 为 K(X) 上由 f 定义的赋值（第 VI 章，§3，第 3 节，例 4）。另一方面，对于 A 的每个本性赋值 w，令 \(\bar{w}\) 为 w 到 K(X) 的扩张，它定义为

\[
\bar {w} \left(\sum_ {j} a _ {j} X ^ {j}\right) = \inf _ {j} (w (a _ {j}))
\]

对于 \(\sum_{j} a_j X^j \in K[X]\)（第 VI 章，§ 10，第 1 节，引理 1）。显然 \(v_f\) 和 \(\bar{w}\) 都是离散且赋范的；对于所有 \(u \in K[X]\)，\(v_f(u) = 0\)（分别地，\(\bar{w}(u) = 0\)）的赋值 \(v_f\)（分别地，\(\bar{w}\)）除有限个外均如此。

为证明命题，只需证明 A{[}X{]} 是赋值 \(v_{f}\) 和 \(\bar{w}\) 的赋值环的交。现在，赋值 \(v_{f}\) 的赋值环的交是 K{[}X{]}。另一方面，对于 \(\sum_{j} a_{j} X^{j} \in K[X]\)，关系 \(\bar{w}\left(\sum_{j} a_{j} X^{j}\right) \geqslant 0\) 等价于“\(w(a_{j}) \geqslant 0\) 对所有 j 都成立”；因此，关系“\(\bar{w}\left(\sum_{j} a_{j} X^{j}\right) \geqslant 0\) 对每个赋值 \(\bar{w}\) 都成立”等价于“\(w(a_{j}) \geqslant 0\) 对所有 j 以及 A 的每个本性赋值 w 都成立”。这就证明了断言。

注。命题 13 的证明中引入的赋值 \(v_f\) 和 \(\bar{w}\) 是 A{[}X{]} 的本性赋值。我们只需证明：若 V 是由赋值 \(v_f\)（\(f\) 不可约）和 \(\bar{w}\)（\(w\) 是本性的）组成的集合，则对每个 \(v' \in V\)，存在元素 \(g \in K(X)\)，它不属于 A{[}X{]}，并且 \(v''(g) \geqslant 0\) 对所有赋值 \(v'' \in V\)（不同于 \(v'\)）都成立；这将证明 \(V - \{v'\}\) 不满足 (AK \(_{\text{II}}\) )，因而结论由第 4 节、命题 6 的推论 2 得出。先假设 \(v'\) 具有 \(\bar{w}\) 的形式：于是可以取 \(g\) 为元素 \(b \in K\)，使得 \(w(b) < 0\)，并且 \(w'(b) \geqslant 0\) 对 A 的本性赋值 \(w'\)（不同于 \(w\)）都成立，因为 \(v_f(b) = 0\) 对每个不可约首一多项式 \(f\)（属于 K{[}X{]}）都成立；满足上述条件的元素 \(b\) 的存在性由第 5 节、命题 9 得出。再假设 \(v'\) 具有 \(v_f\) 的形式，其中 \(f \in K[X]\) 是次数为 \(m\) 的不可约首一多项式；于是可以取 \(g = a/f\)，其中 \(a \in A\)。有 \(v_h(g) \geqslant 0\)（对于每个不可约首一多项式 \(h \neq f\)，它属于 K{[}X{]}）；还需选择 \(a \in A\)，使得对于 A 的每个本性赋值 \(w\)，\(w(a)\) 至少等于元素 \(w(c_i)\) 的最大下界，其中 \(c_i\) 是 \(f (1 \leqslant i \leqslant m)\) 的系数；现在，这样的 \(a \in A\) 的存在性由 (AK \(_{\text{III}}\) ) 及第 5 节、命题 9 得出。

我们还可以说（第 6 节，定理 4），A{[}X{]} 的高度为 1 的素理想是：

\begin{enumerate}
\def\labelenumi{(\arabic{enumi})}
\item
  形如 \(\mathfrak{p}\mathbf{A}[\mathbf{X}]\) 的素理想，其中 \(\mathfrak{p}\) 是 \(\mathbf{A}\) 的高度为 1 的素理想；
\item
  形如 \(\mathfrak{m} \cap \mathrm{A}[\mathrm{X}]\) 的素理想，其中 \(\mathfrak{m}\) 是 \(\mathrm{K}[\mathrm{X}]\) 的（必为主理想的）素理想。
\end{enumerate}

后者的特征是它们与 A 的交化为 0。

